\documentclass[11pt,a4paper]{article}
\usepackage[T1]{fontenc}
\usepackage[utf8]{inputenc}
\usepackage{lmodern}
\usepackage[margin=25mm,headheight=15pt]{geometry}
\usepackage{amsmath,amssymb,amsthm,mathtools}
\usepackage{booktabs,longtable,array,tabularx}
\usepackage{xcolor}
\usepackage{enumitem}
\usepackage{microtype}
\usepackage{listings}
\usepackage{fancyhdr}
\usepackage[hypertexnames=false]{hyperref}
\usepackage{xurl}
\definecolor{navy}{HTML}{16324F}
\definecolor{teal}{HTML}{176B70}
\definecolor{pale}{HTML}{EFF5F7}
\hypersetup{colorlinks=true,linkcolor=navy,urlcolor=teal,citecolor=teal,
 pdftitle={A 239-state Turing machine for ZFC consistency},
 pdfauthor={Research notes prepared for Vladimir Reshetnikov}}
\pagestyle{fancy}
\fancyhf{}
\fancyhead[L]{\small\color{navy}ZFC consistency machines}
\fancyhead[R]{\small Research submission / 19 October 2026}
\fancyfoot[C]{\thepage}
\setlength{\parskip}{5pt}
\setlength{\parindent}{0pt}
\setlength{\emergencystretch}{2em}
\setlist{itemsep=3pt,topsep=5pt,leftmargin=*}
\lstset{basicstyle=\small\ttfamily,breaklines=true,columns=fullflexible,
 keepspaces=true,frame=single,rulecolor=\color{pale},backgroundcolor=\color{pale},
 xleftmargin=4pt,xrightmargin=4pt}
\newtheorem{theorem}{Theorem}[section]
\newtheorem{lemma}[theorem]{Lemma}
\newtheorem{proposition}[theorem]{Proposition}
\newtheorem{corollary}[theorem]{Corollary}
\theoremstyle{definition}
\newtheorem{definition}[theorem]{Definition}
\newtheorem{remark}[theorem]{Remark}
\newcommand{\ZFC}{\mathrm{ZFC}}
\newcommand{\ZF}{\mathrm{ZF}}
\newcommand{\Con}{\operatorname{Con}}
\newcommand{\Halt}{\operatorname{Halt}}
\newcommand{\CI}{\mathsf{CI}}
\newcommand{\CD}{\mathsf{CD}}
\newcommand{\CB}{\mathsf{CB}}
\newcommand{\CC}{\mathsf{CC}}
\newcommand{\CH}{\mathsf{CH}}
\newcommand{\file}[1]{\texttt{\path{#1}}}
\newcommand{\pin}{13ec39ed43148b1807f51ed6e6528e1c84659562}

\begin{document}
\begin{titlepage}
\vspace*{12mm}
{\large\color{teal}COMPUTABILITY / SET THEORY / VERIFIED COMPILATION\par}
\vspace{14mm}
{\Huge\bfseries\color{navy}A 239-state machine\par}
\vspace{5mm}
{\LARGE Nonhalting equivalent to $\Con(\ZFC)$\par}
\vspace{3mm}
{\large Source rotations and a context-dependent state fold\par}
\vspace{9mm}
{\large A research submission for the ProveIt project\par}
\vspace{12mm}
The requested target is a literal binary Turing machine $M$ with fewer than
240 working states and an ordinary proof of
\[
 \ZFC\vdash\bigl(\neg\Halt(M)\leftrightarrow\Con(\ZFC)\bigr).
\]
The report separates exact finite computations, universally quantified
runtime arguments, source-program equivalences, and consistency calibrations.
The companion archive contains the inputs and scripts needed to inspect and
reproduce the reported experiments.
\vfill
\input{status.tex}
\vfill
{\large Prepared for Vladimir Reshetnikov\par}
{19 October 2026, UTC\par}
\vspace{5mm}
{\small Baseline repository commit:\par
\texttt{\pin}\par}
\end{titlepage}

\pagenumbering{roman}
\tableofcontents
\newpage
\pagenumbering{arabic}

\section*{Results at a glance}
\addcontentsline{toc}{section}{Results at a glance}

\textbf{The construction has 239 working states, including blank-tape
initialization.}  An ordinary mathematical proof, composed with the pinned
ProveIt source/theory calibration, establishes
\[
 \ZFC\vdash\bigl(\neg\Halt(M_{239})\leftrightarrow\Con(\ZFC)\bigr).
\]
The new ingredients have independent finite certificates and complete
unbounded runtime checks.  A serialized Lean/Rocq or PA endpoint for this
exact table is not supplied; that distinction concerns formal artifact
status, not an omitted mathematical simulation case.

\begin{center}
\begin{tabularx}{\textwidth}{@{}p{.31\textwidth}X@{}}
\toprule
Artifact & Established result\\
\midrule
Final binary machine & Exactly 239 working states and 478 rows, a
separate halt outcome, and an all-zero two-way infinite initial tape.\\
Source rotation certificate & 225 to 224 distinct source DAG nodes;
maximum path length 51 to 48; maximum zero run remains six.  All
6,455 instructions and their complete control flow are preserved.\\
Final state-fold certificate & Six two-cell block identities and a
complete incoming-edge audit prove the reduction from the repaired
240-state reference for every initial binary tape.\\
Final unbounded runtime & All 8,386 instruction guards pass, with
2,247,513 symbolic rules, arbitrary natural register values and gap,
and an arbitrary nonempty finite spare list.\\
Rejected 239-state patches & Concrete whole-tape counterexamples show
why scanner shortening and uncharged decoder reuse fail.\\
Theory filters & Two precisely weakened container theories have set
models whose existence ZFC proves, excluding the required consistency
calibration if ZFC is consistent.\\
Reusable search & A portable source certificate checker, symbolic
runtime checker, and bounded local state-fold enumerator.\\
\bottomrule
\end{tabularx}
\end{center}

The source improvement uses 25 tree rotations in 13 stages.  The last
two cross an explicit control-flow boundary; their validity follows from
an exact finite graph certificate rather than a generic associativity
assumption.  The initial binary compilation needs a context repair,
returning the count to 240.  Splitting one decoder state's incoming
roles by the bit just written behind the head allows that helper to be
repurposed and two remaining roles to share a state.  This completes
the reduction to 239.

The archive's \file{README.md} and \file{research_status.json} identify
which artifacts are accepted at their stated layer and which are rejected
experiments.  The accepted table is under \file{candidate239/}; earlier
239-state patches retained as diagnostics are not the final machine.

\section{The exact target and the proof obligations}
\label{sec:target}

There is no logical conflict between the requested biconditional and the
second incompleteness theorem.  A theory may prove that a particular machine
does not halt exactly when the theory is consistent, while being unable to
prove either consistency or nonhalting.  The construction must preserve this
distinction throughout the compiler pipeline.

\begin{definition}[Machine convention]
A machine has one tape indexed by $\mathbb Z$, symbols $\{0,1\}$, blank symbol
$0$, a specified start state, and one head initially at position $0$.
Each working state has exactly two instructions.  An instruction writes a
symbol, moves exactly one cell left or right, and enters a working state or
the distinguished terminal outcome.  The terminal outcome is not a working
state.  The initial tape is all zero.  There is no free input word or uncharged
initialization program.
\end{definition}

For a fixed literal table $M$, choose the usual primitive-recursive coding
of finite configurations and transitions, and write
\[
 \Halt(M)\equiv\exists t\,\operatorname{HaltsBy}(\ulcorner M\urcorner,t).
\]
For a fixed effective presentation of set theory, write
\[
 \Con(\ZFC)\equiv
 \forall p\,\neg\operatorname{Prf}_{\ZFC}(p,\ulcorner\bot\urcorner).
\]
Changing reasonable proof codings is harmless only after the effective
translations between their finite proofs have been justified.  A claim about
external consistency alone is not automatically a theorem inside ZFC.

The useful positive route is to establish
\begin{align}
 \ZFC&\vdash \Halt(M)\leftrightarrow
       \exists p\,\operatorname{Prf}_{T}(p,\ulcorner\bot\urcorner),
       \label{eq:machine-calibration}\\
 \ZFC&\vdash \Con(T)\leftrightarrow\Con(\ZFC).
       \label{eq:theory-calibration}
\end{align}
Classical logic then gives the requested nonhalting biconditional.  In a
fully spelled-out implementation, the first line has four interfaces:
the proof calculus, the enumeration of finite candidate proof words, the
register program, and the literal binary machine including startup.
Each interface needs soundness and completeness in the direction relevant
to halting.  A correct instruction interpreter with an incorrect initial
register file is not sufficient.

\begin{center}
\begin{tabularx}{\textwidth}{@{}p{.24\textwidth}X@{}}
\toprule
Obligation & Required evidence\\
\midrule
Theory & Effective presentation of $T$; a ZFC-provable consistency
equivalence, or sufficiently explicit proof/model translations.\\
Calculus & A finite proof of falsity exists in the enumerated syntax exactly
when $T$ is inconsistent.  All side conditions and equality conventions
must be accounted for.\\
Enumerator & Every accepted certificate is sound; a refutation has some
finite encoded candidate that the actual enumeration reaches.\\
Source & Register operations, branch conventions, root rollover, and
initial register state implement the enumerator.\\
Binary runtime & Every source operation is simulated on the entire tape for
unbounded register values and spare lists, with positive finite progress.\\
Literal machine & State census, complete transition table, blank startup,
and all modifications to reused control states are correct.\\
\bottomrule
\end{tabularx}
\end{center}

The user asks for provability in ZFC; a serialized PA theorem is a stronger
engineering goal, not an extra requirement imposed here.  Similarly, an
ordinary mathematical proof need not be a Lean or Rocq kernel theorem.
Nevertheless a Python success record is not itself a kernel-checked proof.
Where the runtime checker is used below, its finite symbolic rules have an
ordinary induction argument; that argument and the remaining interfaces
are stated separately.

\section{Prior art and the pinned 240-state baseline}
\label{sec:baseline}

The user's two prior-art references are the Busy Beaver Challenge logical
independence overview and the $\mathtt{m254}$ development by Mira-Cult and
Jeff Ketchersid.  The latter's README states an end-to-end Coq verification
of the 254-state machine: it halts exactly when its SIR theory is
inconsistent, with the advertised ZF/ZFC consistency calibration.  Andrew
James Wade supplied the earlier compiler framework.  This report does not
rerun that external proof assistant development or transfer its endpoint
to a changed table by assertion.  See the primary sources in
Section~\ref{sec:references}.

The ProveIt baseline examined here is \file{candidate-rotation240}, at the
pinned commit shown on the cover.  The archive includes that directory's table data, source
paths, construction record, attribution, and licenses.  The complete
table and source hashes are recorded in \file{validation/baseline_literal_checks.json}
and the archive manifest.

\begin{center}
\begin{tabular}{@{}lrr@{}}
\toprule
Quantity & Inherited labels & Canonical labels\\
\midrule
Working states & 240 & 240\\
Transition rows & 480 & 480\\
Start state & 95 & 0\\
Dispatch state & 96 & 84\\
Terminal label & 254 & 240\\
\bottomrule
\end{tabular}
\end{center}

\file{validation/check_baseline.py} independently recovers the rooted
transition-graph isomorphism between the two tables and checks every action,
including the final write and movement.  It also checks the supplied
renumbering against that recovered map.  All working states lie in the
rooted transition graph; this is graph reachability and does not assert that
every state is actually visited from the blank tape.

The baseline already contains several transformations.  In particular,
the most recent source rotation changes inherited state 142, and the
240-state fold deletes inherited state 90 and changes the route through
state 72.  Earlier context-dependent folds, including the ones associated
with states 27/15 and 39/17, remain part of the contract.  These are
\emph{prior results}, not new reductions established by this submission.

The current repository README distinguishes ordinary mathematical
arguments from incomplete Lean endpoint integration.  It explicitly leaves
the complete 240-state Lean endpoint and a serialized arithmetic
equivalence open.  Its decoder-coverage material is partial.  Our new
symbolic runtime computation below covers all current instruction sites,
but does not close the missing proof-assistant theorem or silently replace
the source-to-calculus calibration with a finite table test.

\subsection{Identifying the exact calibrated source data}

The portable \file{source_interface/check_source_interface.py} starts from
six unchanged Lean source files fetched at the pinned commit.  It parses
all 203 literal rows and the complete conditional lookup tree in
\texttt{ProgramRouting245}, verifies every selected row and the root 202,
and applies the literal overrides at rows 5, 6, and 169.  Expanding that
actual source matches every one of the baseline's 6,455 ordered
path/instruction/operand triples.  Exactly 202 internal rows are reachable;
row 4 is unused.  This independently identifies the JSON source on which
our source and runtime certificates operate with the program in the
previous consistency calibration.

The same check extracts \texttt{L2.AS}, every branch of \texttt{L2.enc},
and the definitions of \file{LogicalSeed246.seed}, \texttt{registers},
and \texttt{initial}.  It checks the empty path, the nine primary values
below, and the zero tail at all higher register indices.  Each included
Lean file's Git blob identifier is recomputed from its complete bytes.
These are exact data and interface checks, not a new Lean kernel replay.

\subsection{A fresh exact blank-start calculation}

Starting at inherited state 95 on the all-zero tape, the first qualified
\emph{root} frame occurs after exactly 409,691 binary transitions.  A root
has an empty program counter; earlier valid dispatches at non-root program
counters do not qualify.  The verified configuration is:
\begin{center}
\begin{tabularx}{\textwidth}{@{}lX@{}}
\toprule
Field & Value\\
\midrule
State / head & $96$ / $-5$\\
Left marker / header & $-6$ / $176$\\
Logical registers $(N,T,W,A,D,H,X,Y,L)$ & $(3,0,0,0,3,0,0,0,0)$\\
Spare registers & 167 entries, every entry equal to 1\\
Additional header gap & 117, so the header is 181 cells to the right of the head\\
One cells & Exactly 351; every other tape cell is zero\\
\bottomrule
\end{tabularx}
\end{center}
The record stores the entire finite support of the tape, not merely its
register interpretation or the count of ones.  It also records all
transition rows visited before this root.  This is a finite initialization
proof by exact execution.  It says nothing by itself about whether the
subsequent infinite computation halts.

\section{A universal runtime recheck on the literal table}
\label{sec:runtime}

The baseline source has 6,455 leaves: 2,596 increments, 1,931 decrements,
128 breaks, 1,799 continues, and one halt.  Splitting each decrement into
zero and positive cases gives
\[
 6,455+1,931=8,386
\]
instruction guards.  The maximum source depth is 51 and the longest run of
zeros in a source path is six.

\subsection{The quantified tape contract}

For a source path $p$, normalize the dispatch head to position 0.  A valid
frame has one marker at $-1$, the bits of $p$ beginning at 0, zero padding,
one header bit at $64+g$ for $g\geq0$, then unary blocks.  A block for a
register value $v\in\mathbb N$ is
\[
 B(v)=0\,1^{v+1}.
\]
There are nine primary blocks, followed by a first spare block of arbitrary
natural value and an arbitrary finite list of additional spare blocks.
The tape exterior is zero.  The permutation from logical register number
to physical block number is
\[
 (6,2,0,4,3,7,8,5,1).
\]

The successor contract checks the whole tape.  Increment changes the
selected value by one and appends a spare block of value 1.  Positive
decrement reduces the value by one and increases the additional gap by one.
All other primary values and the existing spare list are retained.
Program-counter overflow is charged as the exact displacement of the
marker/head and the corresponding increase in gap.  Halt is the actual
terminal label, not a timeout or an unrecognized runtime frame.

\subsection{Why the computation is not a bounded simulation}

The checker represents each run length by an affine expression in
independent natural variables.  Its variables include the additional gap,
nine primary values, the first spare value, and the total width of the
remaining spare list.  A positive tested register is represented as $u+1$,
while a zero-tested register is exactly zero.  This representation ranges
over all natural inputs admitted by each guard.

Three rules are used:
\begin{enumerate}
\item A literal transition applies one row of the supplied binary table.
\item A self-loop that moves in the scan direction traverses a constant-bit
run.  The rule follows by induction on the natural run length, including
length zero where applicable.
\item A checked one-block traversal is lifted to an arbitrary finite list
of blocks by structural induction on that list.  Its opaque left and right
contexts are retained, so the result is a whole-tape statement.
\end{enumerate}

For the four list-traversal entries, every one-block theorem is rederived
from the literal input table.  Besides $B(v)$, use
\[
 C(v)=1\,0\,1^v.
\]
The relevant entries are 73 rightward $B\to B$, 95 rightward $B\to C$,
48 leftward $B\to B$, and 112 leftward $C\to B$.  Each traversal takes
exactly $v+2$ transitions.  The normalization identity
$C(v)1=1B(v)$ extends to a finite block list by induction, with its reversed
version used when scanning left.  No assertion about arbitrary list
traversal is accepted without these base and step obligations.

Each successful instruction derivation ends at the exact expected frame
with an affine duration whose coefficients are nonnegative and whose
constant term is positive.  Thus one source step has positive finite
binary duration for every instantiated input.  A limit of 4,000 symbolic
proof-construction steps only limits the search for a derivation; it does
not bound register values, spare-list length, gap, or actual binary time.

\subsection{Result and evidence boundary}

The fresh run reports:
\begin{center}
\begin{tabular}{@{}lr@{}}
\toprule
Quantity & Verified result\\
\midrule
Guarded source cases & 8,386 of 8,386\\
Symbolic rule applications & 2,322,575\\
Minimum constant term in the duration & 13 binary transitions\\
Literal transition rows covered & 480 of 480\\
Generic one-block schemas & 4 of 4\\
\bottomrule
\end{tabular}
\end{center}
The report is \file{validation/baseline240_universal.json}; per-case affine
durations and rule counts are in the companion \file{baseline240_universal_cases.json}.
The checker and helper are adapted from the earlier research checker;
\file{validation/adapter_provenance.json} records that relationship.  This
submission does not claim to have invented the symbolic method.

The ordinary proof is the rule-soundness induction just described together
with the successful finite derivations.  The implementation, Python
interpreter, and mathematics of the rule checker remain outside any
proof-assistant trusted kernel.  Source equivalence, theory calibration,
and a formal internalization of the entire checker are separate issues.

\section{Two tempting 239-state tables and their exact failures}
\label{sec:rejected}

The archive preserves rejected candidates as diagnostics.  A concrete
counterexample to an intended runtime contract is enough to reject a
claimed compilation proof.  Such a counterexample need not be reachable
from the blank tape, and it does not prove that the candidate lacks every
possible alternative consistency calibration.

\subsection{Shortening the backward scanner}

The backward scanner's successive zero-counting states include
$4,6,9,13,19,29,41,58$.  Reading a one during the count returns to the
ones scan in state 4.  Reading zero in state 58 enters the returning scan
at state 84.  The eight-zero threshold distinguishes the zero exterior
to the left of the marker from interior program-counter data during
intermediate phases.

Deleting state 58 and changing
\[
 (41,0):0L58\quad\longmapsto\quad0L84
\]
produces a literal table with 239 working states.  It is not a valid
drop-in replacement.  The source's maximum zero run is six, but an
arithmetic selection phase can introduce an additional zero.  The scanner
therefore needs more than a static check of the source leaves alone.

All 8,386 guards were screened with zero values except the tested positive
register, which was set to one, one zero spare block, and no additional
gap.  Exactly 8,287 pass and 99 fail.  The 99 failures involve the 50 leaves
with a six-zero run: 49 decrement leaves, each with two failing guards,
and one increment leaf.

A compact witness is source leaf 40,
\[
 p=\mathtt{11101001001000000},\qquad\CD(1),
\]
with all primary values zero, one zero spare, and $g=0$.  The shortened
table redispatches after 650 transitions with the wrong tape: the symmetric
difference from the expected set of one cells is $\{13,15\}$.  The
unmodified 240-state table satisfies the same contract.  The positive
fixture at this leaf also fails, with differences at
$\{14,15,17,18\}$.  These witnesses are in
\file{scanner/scanner_probe.json}; the script constructs the patch from
the pinned baseline, rather than relying on an unexplained extra table.

\subsection{Reassociating a shared push--copy sequence}

At inherited decoder state 98, the source groups three closed fragments as
$A;(B;C)$, where the fragments are the current push operation, parameter
copy, and the following source fragment.  Existing state 113 supplies the
grouping $(A;B)$.  The obvious reassociation yields the patch
\begin{center}
\begin{tabular}{@{}cll@{}}
\toprule
Row & Baseline & Tentative replacement\\
\midrule
$98/0$ & $0R23$ & $0R113$\\
$98/1$ & $1R74$ & $1R16$\\
$74/0,74/1$ & Present & Deleted\\
\bottomrule
\end{tabular}
\end{center}
There are 13 source occurrences of this reassociation, 6,455 leaves remain,
and the literal table has 239 working states.  The exact blank startup is
unchanged because neither altered state is used before the qualified root.
This same basic source idea was already examined against an earlier
245-state table; its appearance here is not a novelty claim.

The new paths contain six seven-zero runs, each at a decrement of register
1.  That is already enough to invalidate the unmodified scanner contract.
There is a second obstruction: a backward decoder route uses state 98 as
part of a context-dependent fold.  Preserving the source meaning of state
98 when entered normally does not preserve its meaning when re-entered
from the left-moving route through state 218.

The finite guard screen fails on 702 of 8,386 cases.  Of these, 690 are
associated with the changed backward route; the other 12 are the two
guards at each of the six seven-zero decrement leaves.  In particular,
the entry through $237/0$ reaches state 218 with left tag zero; the
entries through $187/1$ and $220/1$ have left tag one and retain the
other route.  The literal patch and all screened failures are in
\file{source_encoding/reassoc239_old.tm} and
\file{source_encoding/reassoc239_finite_screen.json}.

This example identifies a real compiler issue: source sharing, path-length
constraints, and decoder context sharing are coupled.  Neither a lower
source-node count nor correct blank initialization proves a smaller
consistency machine.

\input{source_result.tex}
\input{binary_result.tex}
\input{fold239.tex}
\input{endpoint239.tex}

\section{A sound register invariant and what it does not remove}
\label{sec:abstract}

Some source leaves may be reached only with a zero register, even though
the generic runtime contract supports arbitrary values.  This could permit
removing a branch or using a cheaper macro.  To test that possibility on the original rotation240 baseline,
\file{semantic/abstract_zeros.py} computes a finite abstract interpretation
of the exact source control-flow graph.

For each register, the domain is
$\{\{0\},\mathbb N_{>0},\mathbb N\}$, with an unreachable state for an
entire source location.  Increment yields positive.  A zero decrement
branch keeps zero; a positive decrement yields unknown after subtraction.
The other instructions retain register facts.  Join is componentwise
union.  Start at the exact first root with logical values
$(3,0,0,0,3,0,0,0,0)$.  The worklist converges because the domain and the
source graph are finite.  A separate pass verifies that every generated
edge is contained in the computed destination facts.

By induction on source execution length, every reachable register vector
belongs to the abstract value recorded at its source location.  This is a
sound invariant, although it loses relations between registers and cannot
distinguish different positive values.

The result contains 6,453 reachable source leaves, 5,534 singleton-zero
facts, and exactly two unreachable leaves:
\[
\begin{array}{ccl}
6430 & \mathtt{1111110010} & \CI(6),\\
6431 & \mathtt{1111110011} & \CC(2).
\end{array}
\]
Only one decrement guard is known to be zero; no tested decrement register
is known positive at its entry.  Although this excludes three of the
8,386 generic guards of that baseline, all 480 of its binary transition rows are still used by the
remaining symbolic instruction proofs and the generic block proofs.

The restricted unbounded rerun covers 8,383 guards and uses 2,310,622
symbolic rules.  There is no free literal row under this invariant.  This
does not show that all rows are required by every reachable run, prove
that no state can be merged, or exclude stronger relational invariants.
It rules out the specific strategy of deleting a row because the current
zero/nonzero analysis never needs it.

Finite semantic signatures of closed source fragments were also screened.
Some conditional coincidences occur on sampled zero-register fixtures, but
none survived the requirement that the necessary zero facts hold at every
occurrence.  These heuristic comparisons are kept separate from the
inductive invariant and are not used as equality theorems.

\input{fold_search.tex}

\input{theory_filters.tex}

\input{alternative_formalisms.tex}

\section{A compiler research direction with a precise objective}
\label{sec:direction}

The successful 239-state construction optimizes the \emph{complete binary
compiler contract} while retaining the already calibrated source calculus.
The boundary-aware source result widens the certified search space beyond
syntactically closed sequence fragments.  The final fold then prices and
repairs the decoder contexts explicitly.  The following design principles
describe what worked here and provide an objective for further reductions.

A compiler search state should carry at least four pieces of information:
\begin{enumerate}
\item The source operations and their exact successor indices, including
both decrement guards and root rollover.
\item The binary path language and the maximum zero run required by every
runtime phase, not only by a dispatch snapshot.
\item Entry context for a reused decoder state: head displacement and the
relevant tape bits behind the head, with a checked invariant establishing
those facts after the proposed modification.
\item The literal set of retained transition rows, including arithmetic,
scan, startup, and any repair states.
\end{enumerate}

The objective is the number of distinct working states in the emitted
reachable transition table, subject to its whole-tape contract.  The
number of source DAG nodes is a useful search heuristic.  It is not the
objective itself.  A one-node saving may cost a new context-repair state;
an extra source node may permit deleting two scanner states.  These
possibilities belong in a joint search rather than independent greedy
passes that assume the previous compiler folds remain valid.

\subsection{A checkable candidate acceptance sequence}

The experiments suggest a practical order for future candidates:
\begin{enumerate}
\item Emit a complete table, a state census, exact source paths, and the
full patch or reproducible generator.
\item Check the source control-flow certificate and the literal initial
tape computation.  Name every inherited fold touched by a changed route.
\item Run the inexpensive complete minimum-register guard screen.  A
single mismatch produces a full configuration witness and rejects the
claimed generic contract.
\item Recompute context invariants on the changed table.  Never transfer
an old invariant merely because the source transformation is sound.
\item Re-prove all unbounded instruction and block contracts from the
changed literal rows.  Prove each binary segment has positive finite time.
\item Compose with source calibration and the consistency theorem, and
only then state the requested biconditional for the exact machine hash.
\end{enumerate}

For a new formalism, add an earlier gate: prove the precise consistency
equivalence and design the complete enumerator before quoting a promising
formula length.  A short set-theory sentence can require a long
well-formedness checker, side-condition routine, or proof search language.
Conversely, a longer sentence may reuse source macros already present.

\section{Reproduction and incorporation}
\label{sec:reproduce}

The archive is intended to be copied into a research directory of ProveIt
without changing existing repository files.  It includes TeX sources,
the rendered PDF, Python and C++ experiment sources, pinned baseline
data, selected result records, and attribution.  Rejected 239-state
tables are explicitly labeled as rejected in the research index.

The root \file{README.md} gives exact commands and describes the quick
and full runs.  Python scripts use the standard library.  The universal
runtime checks take several minutes and must be run with assertions
enabled; do not use Python's \texttt{-O} flag.  C++ searches require a
C++17 compiler and are bounded experiments with their stated parameters.
Their absence of a hit is not a state lower bound.

For the report, run \texttt{pdflatex} twice from the archive root.  The
PDF has been rendered to page images and visually inspected.  The package
manifest hashes the included source/data artifacts; timing fields and
environment descriptions in a fresh rerun may naturally differ.  A
successful replay should be judged by semantic statuses, exact counts,
pins, and verified identities, rather than byte equality of timestamps.

The supplied baseline Lean files are source references.  This archive
does not claim a newly checked Lean or Rocq endpoint.  In particular,
the ordinary induction proofs and the source certificates must be
distinguished from a kernel theorem that imports a literal candidate table.

\section{Primary sources and attribution}
\label{sec:references}

\begin{enumerate}[label={[\arabic*]}]
\item Busy Beaver Challenge, \emph{Logical independence}.
\url{https://wiki.bbchallenge.org/wiki/Logical_independence}.
Used as the user's prior-art overview; the original verification repository
is the primary source for machine-specific claims.

\item Mira-Cult; Jeff Ketchersid, \emph{ZF/m254}, branch
\texttt{zf-254}, in Jeff Ketchersid's busy-beaver repository.
\url{https://github.com/jeff87654/busy-beaver/tree/zf-254/ZF/m254}.
Original compact machine/calculus and Coq verification; the README describes
the SIR, ZF, and ZFC consistency interfaces.

\item Vladimir Reshetnikov et al., \emph{ProveIt},
\file{Computability/BusyBeaver/Research/ZFC/candidate-rotation240/README.md},
commit \texttt{\pin}.
\url{https://github.com/VladimirReshetnikov/ProveIt/tree/13ec39ed43148b1807f51ed6e6528e1c84659562/Computability/BusyBeaver/Research/ZFC/candidate-rotation240}.
The precise baseline, construction history, and endpoint limitations.

\item ProveIt incoming research index.
\url{https://github.com/VladimirReshetnikov/ProveIt/tree/13ec39ed43148b1807f51ed6e6528e1c84659562/docs/incoming}.
Earlier experiments and reconciliations must be consulted before assigning
novelty to a local simplification.

\item ProveIt native-theory and compiler research notes at the same commit:
\file{native_theory254.md}, \file{compiler_route.md},
\file{mira_backend.md}, \file{full_grammar_experiment.md},
\file{dual_quantifier_gap.md}, and the references linked in the theory
sections.  Paths are relative to
\file{Computability/BusyBeaver/Research/ZFC} unless a full path is supplied.
\end{enumerate}

Original artifacts retain their upstream license notices in
\file{baseline/rotation240/licenses}.  New experiments and their
adaptation history are identified in \file{ATTRIBUTION.md}.  A new source
search result does not erase attribution for the underlying calculus,
machine, compiler, verification, or earlier source and context folds.

\end{document}
